% Source: img/2010/jiangsu/q7_fig1.png, question 7.
% Complete algorithm flowchart redraw: initialize S and n, accumulate 2^n,
% test S>=33, loop on N, and output S on Y.
\begin{tikzpicture}[
  >=Stealth,
  line width=0.55pt,
  line cap=round,
  line join=round,
  every node/.style={font=\normalsize,anchor=center,align=center,
    execute at begin node={\everypar{}\leftskip=0pt\rightskip=0pt\parindent=0pt
      \hangindent=0pt\hangafter=1\parshape=0}},
  flowbox/.style={draw,minimum height=.70cm,inner xsep=4pt,inner ysep=2pt,
    align=center},
  terminal/.style={flowbox,rounded corners=.18cm,minimum width=1.55cm,
    minimum height=.74cm},
  io/.style={flowbox,shape=trapezium,trapezium left angle=72,
    trapezium right angle=108,minimum height=.80cm},
  decision/.style={flowbox,shape=diamond,aspect=2.85,minimum width=4.10cm,
    minimum height=1.05cm,inner sep=1pt},
  flowarrow/.style={->,line width=0.55pt},
]
  \node[terminal] (start) at (0,11.20) {开始};
  \node[flowbox,minimum width=2.15cm] (initS) at (0,9.80)
    {$S\leftarrow 1$};
  \node[flowbox,minimum width=2.15cm] (initN) at (0,8.45)
    {$n\leftarrow 1$};
  \node[flowbox,minimum width=4.45cm] (accum) at (0,6.90)
    {$S\leftarrow S+2^n$};
  \node[decision] (test) at (0,5.20) {$S\geqslant 33$};
  \node[flowbox,minimum width=2.95cm] (increment) at (4.10,6.90)
    {$n\leftarrow n+1$};
  \node[io,minimum width=2.40cm] (output) at (0,3.40) {输出 $S$};
  \node[terminal] (finish) at (0,2.00) {结束};
  \coordinate (merge) at (0,7.75);

  \draw[flowarrow] (start.south) -- (initS.north);
  \draw[flowarrow] (initS.south) -- (initN.north);
  \draw (initN.south) -- (merge);
  \draw[flowarrow] (merge) -- (accum.north);
  \draw[flowarrow] (accum.south) -- (test.north);

  % N branch: increment n, then return to the accumulation step.
  \draw[flowarrow] (test.east) -- (4.10,5.20) -- (increment.south);
  \draw[flowarrow] (increment.north) -- (4.10,7.75) -- (merge);
  % Keep the branch label clear of the horizontal connector.
  \node[anchor=south west,inner sep=0pt] at (1.78,5.42) {N};

  % Y branch: output and terminate.
  \draw[flowarrow] (test.south) -- (output.north);
  \node[anchor=west] at (0.18,4.38) {Y};
  \draw[flowarrow] (output.south) -- (finish.north);

  \path[use as bounding box] (-2.90,1.62) rectangle (5.25,11.60);
\end{tikzpicture}
